\documentclass[11pt, letterpaper]{article}

% --- Idioma y codificación ---
\usepackage[spanish]{babel}
\usepackage[utf8]{inputenc}
\usepackage[T1]{fontenc}
\usepackage{lmodern}

% --- Matemáticas ---
\usepackage{amsmath, amssymb, amsthm}

% --- Formato ---
\usepackage[margin=2.5cm]{geometry}
\usepackage{parskip}
\usepackage{booktabs}
\usepackage{hyperref}
\hypersetup{colorlinks=true, linkcolor=blue!70!black, urlcolor=blue!70!black}

% --- Entornos de teoremas ---
\theoremstyle{definition}
\newtheorem{definition}{Definición}[section]
\theoremstyle{plain}
\newtheorem{theorem}{Teorema}[section]
\newtheorem{lemma}[theorem]{Lema}
\newtheorem{corollary}[theorem]{Corolario}
\theoremstyle{remark}
\newtheorem{remark}{Observación}[section]

\title{%
  T03 --- Demostración formal:\\[4pt]
  Longitud esperada de cadena $= \alpha$
}
\author{Curso Linux--C--Rust --- B15 Estructuras de datos en C}
\date{}

\begin{document}
\maketitle
\tableofcontents

% ===================================================================
\section{Modelo}
% ===================================================================

\begin{definition}[Tabla hash]
Un array $T[0..m{-}1]$ de $m$ buckets. Cada bucket $T[j]$ contiene
una lista enlazada (cadena). Se insertan $n$ claves
$k_1, k_2, \ldots, k_n$.
\end{definition}

\begin{definition}[Hashing uniforme simple (SUHA)]
Para cada clave $k$, $h(k)$ se distribuye uniformemente en
$\{0, 1, \ldots, m{-}1\}$, e independientemente de las demás claves:
\[
  \Pr[h(k) = j] = \frac{1}{m}
  \quad \forall\, j \in \{0, \ldots, m{-}1\}
\]
\end{definition}

\begin{definition}[Factor de carga]
$\alpha = n/m$.
\end{definition}

\begin{definition}
Sea $L_j$ la longitud de la cadena en el bucket~$j$, es decir, el
número de claves mapeadas a $T[j]$.
\end{definition}

% ===================================================================
\section{Demostración 1:
  \texorpdfstring{$E[L_j] = \alpha$}{E[Lj] = alpha}
  por linealidad de la esperanza}
% ===================================================================

\begin{theorem}\label{thm:expected}
Bajo SUHA, para todo bucket $j \in \{0, \ldots, m{-}1\}$:
\[
  E[L_j] = \frac{n}{m} = \alpha
\]
\end{theorem}

\begin{proof}
Definimos variables indicadoras. Para cada clave $k_i$
($1 \leq i \leq n$) y cada bucket~$j$:
\[
  X_{ij} = \begin{cases}
    1 & \text{si } h(k_i) = j \\
    0 & \text{si no}
  \end{cases}
\]

La longitud de la cadena en el bucket~$j$ es:
\[
  L_j = \sum_{i=1}^{n} X_{ij}
\]

Por linealidad de la esperanza (que \textbf{no} requiere
independencia):
\[
  E[L_j] = \sum_{i=1}^{n} E[X_{ij}]
         = \sum_{i=1}^{n} \Pr[h(k_i) = j]
\]

Por SUHA, $\Pr[h(k_i) = j] = 1/m$ para todo~$i$:
\[
  E[L_j] = \sum_{i=1}^{n} \frac{1}{m} = \frac{n}{m} = \alpha
  \qedhere
\]
\end{proof}

\begin{remark}
El resultado vale para \textbf{todo} bucket~$j$, no solo ``en promedio
sobre los buckets''. Cada bucket individual tiene longitud esperada
exactamente~$\alpha$.
\end{remark}

% ===================================================================
\section{Demostración 2:
  \texorpdfstring{$L_j \sim \mathrm{Binomial}(n, 1/m)$}%
  {Lj ~ Binomial(n, 1/m)}}
% ===================================================================

\begin{theorem}\label{thm:binomial}
Bajo SUHA, $L_j$ sigue una distribución binomial con parámetros $n$ y
$p = 1/m$.
\end{theorem}

\begin{proof}
Las variables $X_{1j}, X_{2j}, \ldots, X_{nj}$ son:
\begin{enumerate}
  \item \textbf{Independientes}: por SUHA, $h(k_i)$ es independiente
    de $h(k_\ell)$ para $i \neq \ell$.
  \item \textbf{Bernoulli con el mismo parámetro}:
    $X_{ij} \sim \mathrm{Bernoulli}(1/m)$ para todo~$i$.
\end{enumerate}

Por definición, la suma de $n$ variables Bernoulli independientes con
el mismo parámetro~$p$ es $\mathrm{Binomial}(n, p)$:
\[
  L_j \sim \mathrm{Binomial}(n, 1/m) \qedhere
\]
\end{proof}

% ===================================================================
\section{Demostración 3: varianza de $L_j$}
% ===================================================================

\begin{theorem}\label{thm:variance}
$\mathrm{Var}[L_j] = \alpha(1 - 1/m)$.
\end{theorem}

\begin{proof}
Como $L_j \sim \mathrm{Binomial}(n, 1/m)$, y la varianza de una
$\mathrm{Binomial}(n, p)$ es $np(1{-}p)$:
\[
  \mathrm{Var}[L_j]
  = n \cdot \frac{1}{m} \cdot \left(1 - \frac{1}{m}\right)
  = \frac{n}{m} \cdot \frac{m-1}{m}
  = \alpha \cdot \frac{m-1}{m} \qedhere
\]
\end{proof}

\begin{corollary}\mbox{}
\begin{enumerate}
  \item \textbf{Desviación estándar}:
    $\sigma_{L_j} = \sqrt{\alpha(1 - 1/m)} \approx \sqrt{\alpha}$
    para $m$ grande.
  \item \textbf{Para $m$ grande} ($m \gg 1$):
    $\mathrm{Var}[L_j] \approx \alpha$, y $L_j$ se aproxima a una
    distribución $\mathrm{Poisson}(\alpha)$.
  \item \textbf{Concentración}: por la desigualdad de Chebyshev,
    \[
      \Pr\!\left[|L_j - \alpha| \geq t\sqrt{\alpha}\right]
      \leq \frac{1}{t^2}
    \]
    Con $\alpha = 1$ y $t = 3$:
    $\Pr[L_j \geq 4] \leq 1/9 \approx 11\%$.
\end{enumerate}
\end{corollary}

% ===================================================================
\section{Demostración 4: la longitud total es exactamente $n$}
% ===================================================================

\begin{theorem}\label{thm:total}
$\displaystyle\sum_{j=0}^{m-1} L_j = n$ con probabilidad~1.
\end{theorem}

\begin{proof}
Cada clave $k_i$ es mapeada a exactamente un bucket (la función hash
es total y determinista para cada clave fija). Por lo tanto, cada
$k_i$ contribuye exactamente~1 a exactamente un~$L_j$:
\[
  \sum_{j=0}^{m-1} L_j
  = \sum_{j=0}^{m-1} \sum_{i=1}^{n} X_{ij}
  = \sum_{i=1}^{n} \underbrace{\sum_{j=0}^{m-1} X_{ij}}_{= 1}
  = n
\]

Verificación de consistencia con el Teorema~\ref{thm:expected}:
\[
  \sum_{j=0}^{m-1} E[L_j]
  = \sum_{j=0}^{m-1} \alpha
  = m \cdot \frac{n}{m} = n \quad \checkmark \qedhere
\]
\end{proof}

% ===================================================================
\section{Demostración 5: longitud máxima esperada}
% ===================================================================

\begin{theorem}\label{thm:max}
Bajo SUHA, con $n = m$ ($\alpha = 1$):
\[
  E\!\left[\max_j L_j\right]
  = \Theta\!\left(\frac{\log n}{\log \log n}\right)
\]
\end{theorem}

\begin{proof}[Prueba (sketch)]
Usamos la cota de la distribución binomial. Para un bucket fijo~$j$:
\[
  \Pr[L_j \geq k]
  = \binom{n}{k} \left(\frac{1}{m}\right)^{\!k}
    \left(1 - \frac{1}{m}\right)^{\!n-k}
  \leq \binom{n}{k} \frac{1}{m^k}
\]

Con $n = m$ y usando $\binom{n}{k} \leq (ne/k)^k$:
\[
  \Pr[L_j \geq k]
  \leq \left(\frac{ne}{km}\right)^{\!k}
  = \left(\frac{e}{k}\right)^{\!k}
\]

\emph{Cota superior.} Sea $k^* = c\log n / \log\log n$ para $c$
suficientemente grande. Entonces $(e/k^*)^{k^*}$ decrece más rápido
que $1/n^2$. Por union bound sobre los $m = n$ buckets:
\[
  \Pr\!\left[\max_j L_j \geq k^*\right]
  \leq n \cdot \frac{1}{n^2} = \frac{1}{n} \to 0
\]

\emph{Cota inferior.} Se muestra (por segundo momento o
inclusión-exclusión) que
$\Pr[\max_j L_j \geq c'\log n / \log\log n] \to 1$ para $c'$
suficientemente pequeño.
\end{proof}

\begin{remark}[Implicación práctica]
Aunque $E[L_j] = 1$ (con $\alpha = 1$), el bucket \textbf{más largo}
tiene $\Theta(\log n / \log\log n)$ elementos. El peor caso de una
búsqueda individual no es $O(1)$, sino
$O(\log n / \log\log n)$.
\end{remark}

% ===================================================================
\section{Verificación numérica}
% ===================================================================

Para $\alpha = 1$ ($n = m$), la distribución $\mathrm{Poisson}(1)$
predice:

\begin{table}[h]
\centering
\begin{tabular}{ccc}
\toprule
Longitud $k$ & $\Pr[L_j = k]$ & $n{=}1000$: buckets con longitud $k$ \\
\midrule
0        & 0.368 & 368 \\
1        & 0.368 & 368 \\
2        & 0.184 & 184 \\
3        & 0.061 & 61  \\
4        & 0.015 & 15  \\
5        & 0.003 & 3   \\
$\geq 6$ & 0.001 & ${\sim}1$ \\
\bottomrule
\end{tabular}
\end{table}

Verificación: $\sum k \cdot \Pr[L_j = k] = 1 = \alpha$. \checkmark

% ===================================================================
\section{Resumen de resultados}
% ===================================================================

\begin{table}[h]
\centering
\begin{tabular}{llp{5.5cm}}
\toprule
\textbf{Resultado} & \textbf{Técnica} & \textbf{Clave} \\
\midrule
$E[L_j] = \alpha$
  & Linealidad de la esperanza
  & $X_{ij} \sim \mathrm{Bernoulli}(1/m)$, sumar $n$ indicadoras \\
$L_j \sim \mathrm{Binomial}(n, 1/m)$
  & Independencia por SUHA
  & Suma de Bernoulli i.i.d. \\
$\mathrm{Var}[L_j] = \alpha(1{-}1/m)$
  & Varianza de la binomial
  & $\approx \alpha$ para $m$ grande \\
$\sum L_j = n$
  & Doble conteo
  & Cada clave en exactamente un bucket \\
$\max L_j = \Theta(\frac{\log n}{\log\log n})$
  & Union bound + segundo momento
  & El peor bucket es superconstante \\
\bottomrule
\end{tabular}
\end{table}

\end{document}
